%% kapa: exam=1 round=19 year=2008 period=1 form=A subject=civil_law subject_ko=민법 seq=7 number=7
%% source: Q-Net 감정평가사 기출문제 (cst00309) artlSeq=1001962 file=extracted/1008/1001962/2008년 제19회 감정평가사 1차시험 기출문제/1교시.hwp
%% conversion: kordoc 4.24.51 -> tools/gen1.py (auto); review=proofread-sonnet; flags=none
\begin{question}{7}
甲은 전원주택의 신축 목적을 乙에게 알리고 그의 토지를 구입하였다. 그러나 그 토지는 관련 법규상 전원주택을 신축할 수 없는 토지였고, 乙은 이 사실을 알고 있었으나 甲에게 알려주지 않았다. 다음 설명 중 옳은 것은? (다툼이 있는 경우에는 판례에 의함)
\begin{choices}
  \item 甲은 동기의 착오를 이유로 매매계약을 취소할 수 없다.
  \item 甲은 사기 또는 동기의 착오를 주장하여 매매계약을 취소할 수 있다.
  \item 甲은 乙에 대하여 불법행위로 인한 손해배상을 청구할 수 없다.
  \item 甲은 사기를 이유로 매매계약을 취소할 수 없다.
  \item 甲은 乙에 대하여 사기에 의한 계약취소나 불법행위로 인한 손해배상청구권을 선택적으로 행사할 수 있을 뿐이다.
\end{choices}
\end{question}
\answer{2}
